% Copyright 2026  Ed Bueler

\documentclass[10pt,
               svgnames,
               hyperref={colorlinks,citecolor=DeepPink4,linkcolor=FireBrick,urlcolor=Maroon},
               usepdftitle=false]{beamer}

\usetheme{Madrid}
\usecolortheme{beaver}
\setbeamercovered{transparent}
\setbeamerfont{frametitle}{size=\large}
\setbeamercolor*{block title}{bg=red!10}
\setbeamercolor*{block body}{bg=red!5}

\usepackage[english]{babel}
\usepackage[latin1]{inputenc}
\usepackage{times}
\usepackage[T1]{fontenc}
% Or whatever. Note that the encoding and the font should match. If T1
% does not look nice, try deleting the line with the fontenc.

\usepackage{empheq,bm,xspace}
%\usepackage{minted}
%\usepackage{hyperref}
\usepackage{tikz}

% If you wish to uncover everything in a step-wise fashion, uncomment
% the following command: 
%\beamerdefaultoverlayspecification{<+->}

\newcommand{\bb}{\mathbf{b}}
\newcommand{\bc}{\mathbf{c}}
\newcommand{\bbf}{\mathbf{f}}
\newcommand{\bl}{\bm{\ell}}
\newcommand{\br}{\mathbf{r}}
\newcommand{\bs}{\mathbf{s}}
\newcommand{\bx}{\mathbf{x}}
\newcommand{\by}{\mathbf{y}}
\newcommand{\bv}{\mathbf{v}}
\newcommand{\bu}{\mathbf{u}}
\newcommand{\bw}{\mathbf{w}}

\newcommand{\bzero}{\bm{0}}

\newcommand{\CC}{\mathbb{C}}
\newcommand{\NN}{\mathbb{N}}
\newcommand{\RR}{\mathbb{R}}
\newcommand{\ZZ}{\mathbb{Z}}

\newcommand{\ddt}[1]{\ensuremath{\frac{\partial #1}{\partial t}}}
\newcommand{\ddx}[1]{\ensuremath{\frac{\partial #1}{\partial x}}}
\renewcommand{\t}[1]{\texttt{#1}}
\newcommand{\Matlab}{\textsc{Matlab}\xspace}
\newcommand{\eps}{\epsilon}

\newcommand{\ds}{\displaystyle}

\newcommand{\twovect}[4]{\ensuremath{{#1}_{#2} =
                            \begin{bmatrix} #3 \\ #4 \end{bmatrix}}}

\newcommand{\rbullet}{{\color{FireBrick} \bullet}}

\newcommand*\circled[1]{\tikz[baseline=(char.base)]{
            \node[shape=circle,draw,inner sep=2pt] (char) {#1};}}

\newcommand{\textbook}{Saxe, \emph{Beginning Functional Analysis}}


\title{Boundary value problems and integral operators}

\subtitle{a calculation for weeks 2 \& 3}

\author{Ed Bueler}

\institute[]{UAF Math 617 Functional Analysis}

\date{Spring 2026}


\begin{document}
\beamertemplatenavigationsymbolsempty

\begin{frame}
  \maketitle
\end{frame}


\begin{frame}{Outline}
  \tableofcontents[hideallsubsections]
\end{frame}

\AtBeginSection[]
{% do not put toc here this time
}


\section{a classic boundary value problem}

\begin{frame}{boundary value problem}

\begin{block}{2-point BVP for the second derivative}
suppose $\alpha,\beta\in\RR$ and $\phi \in C([0,1])$ are given.  solve for $u(x)$:
\begin{align*}
-u''(x) &= \phi(x) \quad \text{ for all } x\in (0,1) \\
u(0) &= \alpha \\
u(1) &= \beta
\end{align*}
\end{block}

\begin{itemize}
\item picture of data $\phi$ (left) and solution $u$ (right):

\vspace{30mm}
\item we will see this is \emph{solvable by hand}, via an integral of the data $\phi(x)$, but this is \emph{a bit hard}, so we warm-up with an easier problem
\end{itemize}
\end{frame}


\begin{frame}{easier: initial value problem}

\begin{block}{IVP for the second derivative\phantom{p}}
suppose $\alpha,\beta\in\RR$ and $\phi \in C([0,1])$ are given.  solve for $u(x)$:
\begin{align*}
-u''(x) &= \phi(x) \quad \text{ for all } x\in (0,1) \\
u(0) &= \alpha \\
u'(0) &= \beta
\end{align*}
\end{block}

\begin{itemize}
\item picture:

\vspace{39mm}
\end{itemize}
\end{frame}


\begin{frame}{solution of IVP}

apply fundamental theorem of calculus twice:

$$-u''(x) = \phi(x)$$

\vspace{70mm}
\end{frame}



\AtBeginSection[] {
    \begin{frame}{Outline}
    \setbeamertemplate{section in toc}[sections numbered]
    \tableofcontents[currentsection]%,hideallsubsections]
    \end{frame}
}


\section{\dots its integral solution}

\begin{frame}{solution of BVP}

apply fundamental theorem of calculus twice?:

$$-u''(x) = \phi(x)$$

\vspace{70mm}
\end{frame}


\begin{frame}{solution of BVP}

\vspace{50mm}

result:
	$$u(x) = \alpha + (\beta - \alpha) x + \int_0^1 k(x,t) \phi(t)\,dt$$
where $k$ is a symmetric \emph{kernel}:
	$$k(x,t) = \begin{cases} t (1 - x), & 0 \le t \le x \le 1 \\
	                         x (1 - t), & 0 \le x \le t \le 1\end{cases}\, = \min\{x,t\} - xt$$
\end{frame}


\begin{frame}{$k(x,t)$ for the BVP}

picture of $k(x,t)$:

\vspace{70mm}
\end{frame}


\begin{frame}{integral solution of BVP}

\begin{itemize}
\item restrict to the $\alpha=\beta=0$ special case
\item recall: $k(x,t) = \min\{x,t\} - xt$
\end{itemize}

\begin{block}{2-point BVP}
the solution of
    $$-u'' = \phi, \quad u(0) = u(1) = 0$$
is
	$$u(x) = \int_0^1 k(x,t) \phi(t)\,dt$$
\end{block}

\begin{itemize}
\item this is a linear \emph{integral operator}
\end{itemize}
\end{frame}


\begin{frame}{the integral operator}

\begin{definition}
given a suitable \emph{kernel} $k(x,t)$, the corresponding \emph{integral operator} is
\begin{gather*}
L \,:\, C([0,1]) \to C([0,1]) \\
(L\phi)(x) = \int_0^1 k(x,t) \phi(t)\,dt
\end{gather*}
\end{definition}

\begin{itemize}
\item $L$ is linear:

\vspace{20mm}
\item is $L$ a matrix?

\vspace{10mm}
\end{itemize}

\end{frame}


\begin{frame}{wait \dots does $\phi\in C([0,1]) \implies L\phi \in C([0,1])$?}

\begin{definition}
given a {\color{FireBrick} suitable} \emph{kernel} $k(x,t)$, the corresponding \emph{integral operator} is
\begin{gather*}
L \,:\, C([0,1]) \to C([0,1]) \\
(L\phi)(x) = \int_0^1 k(x,t) \phi(t)\,dt
\end{gather*}
\end{definition}

\begin{itemize}
\item what does ``suitable'' mean?
\item for $\phi \in C([0,1])$, we see that $(L\phi)(x)$ is well-defined if $k(x,t) \phi(t)$ is (Riemann) integrable in $t$ for every $x$
\item what conditions on $k(x,t)$ would assure that $\phi\in C([0,1])$ implies $L\phi \in C([0,1])$?
  \begin{itemize}
  \item[$\circ$] does $k(x,t)$ need to be continuous?
  \end{itemize}
\end{itemize}
\end{frame}


\begin{frame}{the antiderivative is a well-behaved integral operator}

\begin{definition}
an \emph{antiderivative operator} is, for example,
\begin{gather*}
A \,:\, C([0,1]) \to C([0,1]) \\
(A\phi)(x) = \int_0^x \phi(t)\,dt
\end{gather*}
\end{definition}

\begin{itemize}
\item the FTC shows this is the antiderivative:

\vspace{10mm}
\item note that $(B\phi)(x) = \int_{0.4}^x \phi(t)\,dt$ is another antiderivative operator
\item these are also examples of \emph{Volterra integral operator} \dots look it up?
\end{itemize}
\end{frame}


\begin{frame}{kernel of the antiderivative operator}

\begin{itemize}
\item FIXME antiderivative example, noting $k(x,t)$ is discontinuous
\end{itemize}
\end{frame}


\begin{frame}{Lipschitz kernels give well-behaved integral operator}

\begin{itemize}
\item FIXME show that if there exists $C\ge 0$ so that $|k(x,t)-k(y,t)|\le C |x-y|$ for  $x,t$ then $L\phi$ is continuous
\end{itemize}
\end{frame}


\begin{frame}{returning to the BVP \dots}

\begin{itemize}
\item FIXME thus $u(x)$ solving BVP is indeed continous
\end{itemize}
\end{frame}


\begin{frame}{bounded kernel is not enough}

\begin{itemize}
\item FIXME show by example a discontinuous $k(x,t)$ so that $L\phi$ is not continuous
\item namely $k(x,t)=1$ if $x\le 0.5$ and $k(x,t)=0$ if $x>0.5$
\end{itemize}
\end{frame}


\section{$C([0,1])$ as a normed vector space, with $\|\cdot\|_\infty$}

\begin{frame}{the supremum norm on $C([0,1])$}

\begin{itemize}
\item FIXME define it
\item there is more on this in Chapters 1 and 2 of \textbook
\end{itemize}
\end{frame}

\begin{frame}{compact sets}

\begin{itemize}
\item FIXME define compact by open cover
\end{itemize}
\end{frame}

\begin{frame}{the unit ball in $C([0,1])$}

\begin{itemize}
\item FIXME unit ball $B$
\item it is not compact
\end{itemize}
\end{frame}


\section{compactness and integral operators}

\begin{frame}{bounded operators}

\begin{itemize}
\item FIXME define $L$ is bounded
\end{itemize}
\end{frame}


\begin{frame}{bounded integral operators}

\begin{itemize}
\item FIXME $\sup_{x} \int_0^1 |k(x,t)|\,dt<\infty$ implies operator $L$ bounded
\end{itemize}
\end{frame}


\begin{frame}{new idea: some integral operators are compact}

\begin{itemize}
\item FIXME define $L$ is bounded
\item FIXME $LB$ is compact because equicontinuous and Arzela-Ascoli
\item more on this in Chapter 2 of \textbook
\end{itemize}
\end{frame}


\section{other views of the boundary value problem}

\begin{frame}{x}

\begin{itemize}
\item FIXME:  Dirac delta, Fourier, well-posedness
\item picture:

\vspace{30mm}
\end{itemize}
\end{frame}


\begin{frame}{approximations}

\begin{itemize}
\item FIXME finite elements
\item picture:

\vspace{30mm}
\end{itemize}
\end{frame}



\begin{frame}{theory in weeks 2 and 3}

\begin{block}{to define:}
\begin{itemize}
\item open and closed sets in metric spaces
\item compact sets in metric spaces
\item equicontinuity of subsets of $C([a,b])$
\item complete metric spaces
\item Hilbert and Banach spaces
\end{itemize}
\end{block}

\begin{block}{to prove:}
\begin{itemize}
\item Heine-Borel theorem
\item Arzela-Ascoli theorem
\item $C([a,b])$ with $\|\cdot\|_\infty$ is complete (thus a Banach space)
\end{itemize}
\end{block}

\begin{itemize}
\item this is Chapter 2 in \textbook
\item note that Assignment 2 posted on \href{https://bueler.github.io/fa/index.html}{\texttt{bueler.github.io/fa}}
\end{itemize}
\end{frame}


\end{document}
